\section{Prescribing the local reduction type}\label{sec:local}

The degeneration at the origin allows us to arrange the same local behavior
for every character summand of the covering curve.  We first explain why the
reduction type of an abelian variety is controlled, locally in a
characteristic-zero family, by a sufficiently large finite torsion module.
We then apply this observation to the Jacobian family and translate its
reduction into the local conditions on the two-dimensional representations.

\subsection{Finite torsion and reduction in a local family}

Throughout this section, a local field is a finite extension of
$\Q_\ell$.  In this subsection, the auxiliary torsion prime $d$ is different
from $\ell$.  If $A$ is an abelian variety, write
$V_d(A)=T_d(A)\otimes_{\Z_d}\Q_d$ for its rational Tate module.

\begin{lemma}[Local constancy of finite torsion]\label{loc:torsion}
Let $L$ be a local field, let $X$ be a smooth $L$-variety, and let
$\mathcal A\to X$ be an abelian scheme.  For $x_0\in X(L)$ and an integer
$N\geq 1$, there is an analytic neighborhood $\mathcal N$ of $x_0$ in
$X(L)$ such that
\[
 \mathcal A_x[N](\overline L)\simeq
 \mathcal A_{x_0}[N](\overline L)
 \qquad (x\in\mathcal N)
\]
as $G_L$-modules.
\end{lemma}

\begin{proof}
The group scheme $\mathcal A[N]$ is finite \'{e}tale over $X$, since the
base has characteristic zero.  Equivalently, the scheme of full level-$N$
structures is a finite \'{e}tale torsor over $X$ under the constant group
$\GL_{2g}(\Z/N\Z)$, where $g$ is the relative dimension.
Choose a finite Galois extension $L'/L$ over which the fiber of this torsor
at $x_0$ splits.  The inverse function theorem for an \'{e}tale morphism gives
a local section through each point of the fiber over a neighborhood of
$x_0$ in $X(L')$.  Shrinking that neighborhood, these sections enumerate
all points in every nearby fiber.  Intersecting its finitely many Galois
conjugates, and shrinking once more if necessary, makes the neighborhood
Galois stable and preserves the Galois permutations of the sections.
The right action of the finite level group likewise preserves the same
permutations after shrinking.  Restricting to $X(L)$ therefore identifies
the fibers as Galois sets with their level-group action, and hence
identifies the underlying torsion modules.
\end{proof}

The finite torsion module will first ensure semistability and then
identify a Frobenius polynomial.  For the second step we use the
polynomial on the \emph{full} Tate module: this retains both the toric and
the lattice pieces of the semistable representation.

\begin{lemma}[Finitely many semistable Frobenius polynomials]
\label{loc:frobenius}
Fix a local field $L$ with residue field of cardinality $Q$, an integer
$g\geq 1$, and a prime $d$ different from the residue characteristic.
For a semistable abelian variety $A/L$ of dimension $g$, let
\[
 P_A(T)=\det\bigl(T-\operatorname{Frob}_L\mid V_d(A)\bigr),
\]
where $\operatorname{Frob}_L$ is any lift of arithmetic Frobenius.
Then $P_A(T)$ belongs to $\Z[T]$, is independent of the lift and of $d$,
and ranges over a finite set depending only on $g$ and $Q$.
Moreover, $P_A$ determines the toric rank of $A$; if that rank is zero,
it also determines the Newton polygon of its good reduction.
\end{lemma}

\begin{proof}
Let $r$ be the toric rank of $A$.  The semistable filtration of
Grothendieck identifies the graded pieces of $V_d(A)$ with the rational
Tate module of a torus of rank $r$, the rational Tate module of an abelian
variety of dimension $g-r$ with good reduction, and a lattice of rank $r$
\cite[Expos\'e IX]{SGA7I}.  The character lattice of the torus and the lattice quotient have finite
unramified Galois action.  Consequently the arithmetic Frobenius eigenvalues on the toric and
lattice graded pieces are respectively $Q$ times roots of unity and
roots of unity.  The middle
piece has the $2(g-r)$ Frobenius eigenvalues of an abelian variety over
$\F_Q$, all of complex absolute value $Q^{1/2}$.

Each graded piece has an integral characteristic polynomial independent
of $d$: for the toric and lattice pieces this follows from the finite-order
integral actions just described, together with the toric Tate twist;
for the abelian piece it is the
usual finite-field Frobenius polynomial.  Inertia acts trivially on the
graded pieces, so the product of these three polynomials is also
independent of the lift of Frobenius.  Every root has complex absolute
value at most $Q$.  Thus the coefficient of $T^{2g-j}$ has absolute value
at most $\binom{2g}{j}Q^j$.  Integral polynomials of fixed degree with
these bounds form a finite set.

Exactly $r$ roots have absolute value $1$, exactly $r$ have absolute value
$Q$, and the others have absolute value $Q^{1/2}$.  Since $Q>1$, these
three absolute values distinguish $r$.  If $r=0$, semistability is good
reduction, and the normalized $\ell$-adic valuations of the roots of
$P_A$ are its Newton slopes.
\end{proof}

\begin{proposition}[Local constancy of the semistable type]
\label{loc:constancy}
Let $L_0$ be a local field, let $X$ be a smooth $L_0$-variety, and let
$\mathcal A\to X$ be an abelian scheme of relative dimension $g$.
For $x_0\in X(L_0)$, there exist a finite extension $L/L_0$ and an analytic
neighborhood $\mathcal N\subset X(L_0)$ of $x_0$ with the following
properties.  For every $x\in\mathcal N$, the abelian variety
$\mathcal A_x$ is semistable over $L$ and has the same toric rank as
$\mathcal A_{x_0}$ over $L$.  If $\mathcal A_{x_0}$ has good reduction
over $L$, then $\mathcal A_x$ has the same Newton polygon over $L$.

In particular, if $\mathcal A_{x_0}\simeq E_0^g$ for an elliptic curve
$E_0/L_0$, every sufficiently nearby fiber has, after the same finite
extension, the same one of the following three properties as $E_0^g$:
\begin{enumerate}
\item good ordinary reduction;
\item good reduction with all Newton slopes equal to $1/2$;
\item semistable reduction of toric rank $g$.
\end{enumerate}
\end{proposition}

\begin{proof}
Choose an odd prime $d$ different from the residue characteristic.
By semistable reduction, we may choose a finite extension $L/L_0$ over
which $\mathcal A_{x_0}$ is semistable and all its $d$-torsion points are
rational.  Lemma~\ref{loc:torsion}, applied after base change to $L$,
shows that, in a neighborhood of $x_0$, all nearby fibers also have
rational $d$-torsion over $L$.

These fibers are semistable over $L$.  Indeed, the matrix of any inertia
element on their $d$-adic Tate module has the form $1+dB$ with $B$
integral.  Each eigenvalue $\lambda$ therefore satisfies
$v_d(\lambda-1)\geq 1$, where $v_d(d)=1$.  The local monodromy theorem
makes $\lambda$ a root of unity.  For odd $d$ the only root of unity
satisfying this inequality is $1$: a root with nontrivial prime-to-$d$ part has
$v_d(\lambda-1)=0$, and a primitive $d^a$th root has valuation
$1/(d^{a-1}(d-1))<1$.  Thus every inertia element acts unipotently, and
Grothendieck's criterion gives semistability
\cite[Expos\'e IX]{SGA7I}.

Let $\mathcal P$ be the finite set of polynomials in
Lemma~\ref{loc:frobenius} for dimension $g$ and the residue cardinality
of $L$.  Choose $j$ so large that reduction modulo $d^j$ is injective on
$\mathcal P$.  A second application of Lemma~\ref{loc:torsion} gives,
throughout a smaller neighborhood,
\[
 \mathcal A_x[d^j](\overline L)\simeq
 \mathcal A_{x_0}[d^j](\overline L)
 \quad\text{as }G_L\text{-modules}.
\]
Their full Frobenius polynomials are congruent modulo $d^j$, hence equal.
Lemma~\ref{loc:frobenius} now gives the assertions about toric rank and
Newton polygon.  The inverse image of this neighborhood under
$X(L_0)\to X(L)$ is a neighborhood of $x_0$, as required.

For the final assertion, choose $L$ so that $E_0$ has good or split
multiplicative reduction.  In the good case an elliptic curve is either
ordinary or supersingular; in the latter case both its Newton slopes
are $1/2$.  In the multiplicative case $E_0^g$ has toric rank $g$.
Enlarging $L$ to rationalize its $d$-torsion preserves these properties,
so the preceding argument applies.
\end{proof}

\subsection{Monodromy of a totally toric abelian variety}

The third possibility in Proposition~\ref{loc:constancy} must force
nonzero monodromy on \emph{each} character summand, including when the
coefficient prime is the residue characteristic.  The rank statement
needed for this is a consequence of the valuation pairing in toric
uniformization.  We give the argument at the residue characteristic as
well, in terms of Kummer extensions.

\begin{lemma}\label{loc:toric-monodromy}
Let $L/\Q_\ell$ be finite and let $A/L$ be an abelian variety of
dimension $g$.  Suppose that $A$ acquires semistable
reduction of toric rank $g$ after a finite extension.  For every prime
$d$, the monodromy operator on the Weil--Deligne representation of
$H^1_{\mathrm{et}}(A_{\overline L},\Q_d)$ has rank $g$.
For $d=\ell$, the Weil--Deligne representation is the one attached to
the potentially semistable representation by $\ell$-adic Hodge theory.
\end{lemma}

\begin{proof}
Monodromy rank does not change after finite extension.  We may therefore
assume that $A$ has split toric uniformization
\[
 A^{\mathrm{an}}\simeq (\Gm^g)^{\mathrm{an}}/\Lambda,
\]
where $\Lambda\simeq\Z^g$ is rational over $L$.  Choosing bases of the
torus and lattice writes the periods as a matrix
$(q_{ij})\in(L^\times)^{g\times g}$.  The valuation pairing is
nondegenerate, so
\begin{equation}\label{loc:valuation-matrix}
 \bigl(v_L(q_{ij})\bigr)_{1\leq i,j\leq g}
 \quad\text{is nonsingular.}
\end{equation}
This is the toric case of the monodromy pairing for a semistable
abelian variety \cite[Expos\'e IX]{SGA7I}; its realization on de Rham
cohomology is described in \cite{ColemanIovita}.

If $d\ne\ell$, the same pairing describes the monodromy map from the
lattice quotient to the toric subspace of the Tate module.  Its
nondegeneracy gives rank $g$.  We now prove this rank assertion for
$d=\ell$.

Toric uniformization supplies an exact sequence
\begin{equation}\label{loc:tate-extension}
 0\longrightarrow\Q_\ell(1)^g\longrightarrow V_\ell(A)
   \longrightarrow\Q_\ell^g\longrightarrow0,
\end{equation}
whose extension classes are the Kummer classes of the $q_{ij}$.
The representation is semistable, by the semistable comparison theorem
for abelian varieties \cite{ColemanIovita}.  Let $L_{\mathrm{ur}}$ be the
maximal unramified subfield of $L$, and write
$D_{\mathrm{st}}$ for Fontaine's semistable functor over $L$.
On the two outer terms of~\eqref{loc:tate-extension}, monodromy is zero.
Consequently it induces a map
\[
 \nu:D_{\mathrm{st}}(\Q_\ell^g)
       \longrightarrow D_{\mathrm{st}}(\Q_\ell(1)^g).
\]
After the standard identifications, Frobenius on these two spaces is
respectively $\sigma$ and $\ell^{-1}\sigma$, with $\sigma$ the Frobenius
of $L_{\mathrm{ur}}/\Q_\ell$.  The identity $N\varphi=\ell\varphi N$
therefore makes $\ker\nu\subset L_{\mathrm{ur}}^g$ stable under $\sigma$.

Suppose that $\nu$ has rank less than $g$.  Finite Galois descent for
$L_{\mathrm{ur}}/\Q_\ell$ then gives a nonzero vector
$a=(a_j)\in\Q_\ell^g\cap\ker\nu$.
Pull~\eqref{loc:tate-extension} back along the line spanned by $a$ in
its trivial quotient.  The resulting extension of $\Q_\ell$ by
$\Q_\ell(1)^g$ is semistable and has zero monodromy, hence is crystalline.
Under Kummer theory the Bloch--Kato crystalline subspace
$H_f^1(L,\Q_\ell(1))$ is the kernel of the valuation map
\cite[Example~3.9]{BlochKato}.  Applying this to each coordinate of the
pulled-back extension gives
\[
 \sum_{j=1}^g a_jv_L(q_{ij})=0
 \qquad(1\leq i\leq g),
\]
contrary to~\eqref{loc:valuation-matrix}.
Thus $N$ has rank $g$ on $V_\ell(A)$.  Dualizing gives the same rank
on $H^1_{\mathrm{et}}(A_{\overline L},\Q_\ell)$.
\end{proof}

\subsection{Open conditions on the original parameter}

Return to the elliptic curve $E/K$, the odd integer $n=3p$, and
$F=K(\mu_n)$.  The Jacobian of the constant-$E$ cover $C_n$ extends over
a neighborhood of $O$ in the parameter curve $E$, with fiber $E^n$ at
$O$, by Proposition~\ref{geo:compact-type}.  Proposition~\ref{loc:constancy}
therefore applies at every finite completion of $F$.  The following statement
expresses the resulting neighborhoods in the original parameter
$t\in U\subset\A^1_K$, so they can later be combined with Hilbert
irreducibility.

\begin{proposition}\label{loc:opens}
For every place $v_0\mid 3p$ of $K$, there is a nonempty open subset
$\Omega_{v_0}\subset U(K_{v_0})$ with the following property.
Suppose that $t\in U(K)$ belongs to $\Omega_{v_0}$ for all $v_0\mid3p$,
and set $J_t=\Jac(C_{n,t})$.  At each place $v\mid3p$ of $F$,
$J_t$ has the same one of the following potential reduction types
as $E^n$:
\begin{enumerate}
\item good ordinary reduction;
\item good reduction with all Newton slopes $1/2$;
\item semistable reduction of toric rank $n$.
\end{enumerate}
For each $v$, the extensions giving these properties may be chosen
in advance, independently of $t$ in the indicated open subsets.
\end{proposition}

\begin{proof}
Fix $v_0\mid3p$.  For each $v\mid v_0$, apply
Proposition~\ref{loc:constancy} to the extended Jacobian family over
$F_v$, at the point $O$ of its parameter curve.  This gives a
neighborhood of $O$ in $E(F_v)$ and a fixed finite extension of $F_v$
on which the asserted reduction type holds.  The inverse images of
these finitely many neighborhoods in $E(K_{v_0})$ have an intersection
that is again a neighborhood of $O$.

Choose $R_0\ne O$ in this intersection, sufficiently close to $O$ and
away from the finitely many excluded parameters of the cover family.
We may also avoid the ramification points of
$x:E\to\mathbf P^1$, since these form a finite set.  The local inverse
function theorem gives a nonempty neighborhood
$\Omega_{v_0}\subset U(K_{v_0})$ of $x(R_0)$ and an analytic lift
\[
 R(t)=(t,y(t))\in E(K_{v_0}),\qquad y(t)^2=h(t),
\]
whose image remains in the chosen intersection.  Shrink the neighborhood
so that $y(t)\ne0$ and it avoids all excluded parameters.

The map
\[
 B_t\longrightarrow E,\qquad (x,y)\longmapsto(x,y(t)y)
\]
is an isomorphism over $K_{v_0}$ carrying the marked point $(t,1)$ to
$R(t)$.  It carries the defining function for the cover to the function
in the constant-$E$ construction: their divisors agree, and both have
value $1$ at $O$.  Thus it identifies the covers and their Jacobians at
every $v\mid v_0$.  The reduction conclusions follow from the choice of
the neighborhoods.  Carrying this out independently at the finitely
many $v_0\mid3p$ proves the proposition.
\end{proof}

\subsection{The local representations of the character summands}

For a prime $\ell$ and a rotation character $\theta$, write
$V_{\theta,\ell}$ for the $\ell$-adic realization of the rank-two system
in Section~\ref{sec:systems}.  Recall from Section~\ref{sec:conventions} that the ordinary
filtration of weight zero has an unramified character on its line and
$\eps_\ell^{-1}$ times an unramified character on its quotient, after
a finite extension.  In the ordinary case below the distinguished
line is already preserved before this extension.

\begin{corollary}\label{loc:characters}
Choose $t$ satisfying the open conditions of
Proposition~\ref{loc:opens}.  Let $v\mid3p$ be a place of $F$, of residue
characteristic $\ell$.  Then, simultaneously for every rotation
character $\theta$, the following conclusions hold.
\begin{enumerate}
\item If $E$ has potentially good ordinary reduction at $v$, then
$V_{\theta,\ell}|_{G_{F_v}}$ is potentially crystalline and potentially
ordinary of weight $0$.
\item If $E$ has potentially good supersingular reduction at $v$, then
$V_{\theta,\ell}|_{G_{F_v}}$ is potentially crystalline and is not
potentially ordinary of weight $0$.
\item If $E$ has potentially multiplicative reduction at $v$, then the
monodromy of $V_{\theta,d}|_{G_{F_v}}$ has rank one for every coefficient
prime $d$, including $d=\ell$.  In particular,
$V_{\theta,\ell}|_{G_{F_v}}$ is not potentially crystalline.
\end{enumerate}
\end{corollary}

\begin{proof}
Choose a finite extension $L/F_v$ on which $J_t$ has the reduction given
by Proposition~\ref{loc:opens}.  In the two good-reduction cases,
crystalline comparison makes $H^1_{\mathrm{et}}(J_{t,\overline L},\Q_\ell)$,
and hence each character summand, crystalline over $L$.

Suppose first that the reduction is ordinary.  The connected--\'{e}tale
sequence of the ordinary $\ell$-divisible group, dualized to $H^1$, gives
a filtration whose subrepresentation is unramified and whose quotient
is $\eps_\ell^{-1}$ times an unramified representation.  This filtration
is functorial in endomorphisms and therefore stable under the rotations.
It restricts to such a filtration on each $V_{\theta,\ell}$.
By Proposition~\ref{sys:properties}, this summand has dimension two
and labeled Hodge--Tate weights $0,1$.  Its unramified part contributes
only weight $0$, whereas the quotient contributes only weight $1$;
thus both parts have dimension one.  This is the required ordinary
filtration.  After replacing $L$ by a finite Galois extension, its
weight-zero line is uniquely determined by its inertia action, since
$\eps_\ell$ has infinite image on the inertia of every finite extension.
Conjugation by $G_{F_v}$ preserves that line, so the line descends to
the original decomposition group.

If the reduction has all slopes $1/2$, the crystalline module of every
summand has the same slope.  These normalized slopes remain $1/2$
after further finite extension.  An ordinary filtration of weight $0$
would provide a crystalline rank-one subrepresentation of parallel
Hodge--Tate weight zero.  For that character weak admissibility forces
Newton slope zero, contradicting the slopes of the summand.  Thus no
such potential ordinary filtration exists.

Finally, suppose that $J_t$ has potentially toric rank $n$.
Lemma~\ref{loc:toric-monodromy} gives rank $n$ for monodromy on its
$2n$-dimensional first cohomology, at every coefficient prime.
Monodromy commutes with the rotations, so preserves the $n$ character
summands.  Each summand has dimension two.  On first cohomology of a
semistable abelian variety one has $N^2=0$, so the rank on each summand
is at most one.  Their ranks sum to $n$ and there are exactly $n$
summands; every rank is therefore one.  Nonzero monodromy persists
after finite extension and excludes potential crystallinity.
\end{proof}
